\SourceLeaf{572}Beobachtungen geführten Rechnung enthält folgende Tafel in Secunden ausgedrückt:
\Needspace{19\baselineskip}
\begin{center}\small
\textit{Mittleres Schwanken der magnetischen Declination während der drei Jahre}\par
1834---1837.\par\smallskip
\begin{tabular}{l|r|r|r||r||r|r|r||r}
&\multicolumn{4}{c||}{8 Uhr Vormittag}&\multicolumn{4}{c}{1 Uhr Nachmittag}\\\cline{2-9}
&I&II&III&Mittel&I&II&III&Mittel\\\hline
April&74&126&205&147&129&101&264&180\\
Mai&192&124&277&207&158&183&210&185\\
Junius&172&171&199&181&95&151&217&162\\
Julius&213&243&287&250&119&184&252&193\\
August&264&253&269&262&175&165&307&225\\
September&162&325&207&241&172&143&161&159\\
October&116&296&216&222&182&202&242&210\\
November&79&205&308&218&170&173&126&158\\
December&132&324&71&206&184&206&154&182\\
Januar&146&274&138&196&174&212&154&181\\
Februar&116&146&164&143&178&183&129&165\\
März&100&109&366&228&127&153&246&183\\\hline
Mittel&157&229&238&211&156&174&213&183
\end{tabular}
\end{center}

Von den einzelnen Beobachtungen mögen hier noch die \textit{grössten} Schwan\-kungen angeführt werden, die im Laufe der drei Jahre bei den vormittägigen und nachmittägigen Declinationen vorgekommen sind. Jene war am 8. October 1835 um $20'1''$ grösser als am 7. October, und die nachmittägige Declination am 24. April 1836 um $13'0''$ grösser als am vorhergehenden Tage. Dagegen ist auch völlige Gleichheit der vormittägigen oder der nachmittägigen Declination an zweien auf einander folgenden Tagen öfters vorgekommen. In den monatlichen Mittelschwankungen rücken natürlich diese Extreme viel näher zusammen; gleichwohl bleibt die grosse Ungleichheit der einzelnen Monate in dieser Beziehung sehr bemerkenswerth, da nach obiger Übersicht das mittlere Schwanken bei der Vormittagsdeclination im März 1837 die Grösse von $6'6''$ hatte, im December 1836 hingegen nur $1'11''$ betrug.

Ob im Allgemeinen zu einer Tageszeit grössere Schwankungen vorherrschen als zu einer andern, ist aus den Resultaten für unsere beiden Stunden mit Sicherheit noch nicht zu entscheiden. Im Mittelwerth stehen im ersten Jahre beide
\SourceLeaf{573}nahe gleich, in den beiden andern überwiegen die Vormittagsschwankungen, aber der Unterschied der Endresultate aus allen drei Jahren $3'31''$ und $3'3''$ ist zu klein, als dass man ihn durch so wenige Jahre für festgestellt halten dürfte, wiewohl in den Mittelzahlen für die einzelnen Monate in der vierten und achten Columne zehn Monate eine Differenz in demselben Sinn gegeben haben.

Wirft man Vormittags- und Nachmittagsbeobachtungen zusammen, so erhält man folgende mittlere Schwankungen:
\begin{center}\small
\begin{tabular}{l|r|r|r||r}
&Jahr I&Jahr II&Jahr III&Mittel\\\hline
April&108&114&237&164\\
Mai&176&156&245&196\\
Junius&139&161&208&172\\
Julius&173&215&270&223\\
August&224&214&289&244\\
September&167&251&185&204\\
October&152&254&229&216\\
November&133&190&235&191\\
December&160&271&120&195\\
Januar&160&245&146&189\\
Februar&150&166&148&155\\
März&114&133&312&206
\end{tabular}
\end{center}
\Needspace{7\baselineskip}
\begin{center}\small
\textit{Mittelwerthe.}\par\smallskip
\begin{tabular}{l|r|r|r||r}
&Jahr I&Jahr II&Jahr III&Mittel\\\hline
Julius---December&170&234&228&213\\
Übrige Monate&143&167&223&181\\
Ganzes Jahr&158&204&226&198
\end{tabular}
\end{center}

Nach den Zahlen der vierten Columne herrschen in den Monaten Julius---December etwas grössere Schwankungen vor, als in den sechs übrigen, aber die Mittelwerthe $3'33''$ und $3'1''$ sind doch wohl zu wenig verschieden, um daraus mit Sicherheit schliessen zu können, dass jene Jahreszeit grössere Schwankungen mehr begünstigt, zumal da der Unterschied nur hauptsächlich in dem einen Jahre 1835---1836 auf diese Art stark hervorgetreten ist.

Sehr kenntlich ist hingegen die Ungleichheit der Veränderlichkeit in den einzelnen drei Jahren gegen einander gehalten; der Mittelwerth für das dritte Jahr ist fast um die Hälfte grösser, als der Mittelwerth für das erste. Das Ge\SourceLeaf{574}neralmittel aus sämmtlichen bisherigen Beobachtungen $3'18''$ könnte daher nach längerer Fortsetzung wohl noch erhebliche Abänderung erhalten.

Dies sind die Resultate, die sich aus den bisherigen täglichen Aufzeichnungen der magnetischen Declination ziehen lassen. Es ist sehr zu wünschen, dass ähnliche Arbeiten an mehrern Orten ausgeführt werden, und an einigen ist seit kurzem schon der Anfang damit gemacht. Wenn, wie in Mailand geschieht, die Beobachtungen nicht nach der Ortszeit, sondern genau gleichzeitig mit den hiesigen angestellt werden, so bietet die Vergleichung der einzelnen Tage noch zu andern Combinationen Gelegenheit dar, welche. wenn sie erst eine etwas beträchtliche Zeit umfassen können, von grossem Interesse sein werden. Die Beobachter, welche es auf eine ähnliche Weise halten, d.\ i.\ ihre Aufzeichnungen zu solchen Zeiten machen, welche mit den hiesigen übereinstimmen, werden daher ersucht, die Resultate aller Tage einzeln mitzutheilen, wobei es jedoch zureicht, sie nur nach Scalentheilen anzugeben, so dass die Verwandlung in Bogentheile erspart werden kann, wenn nur zugleich die nöthigen Reductionselemente bemerkt werden.

